\section{概念介绍}\label{Concepts}
此处对一些概念做简要介绍, 可以作为选读内容.
\subsection{编辑器 \& 编译器 \& IDE}
编辑器指文本编辑器, 常见的如微软的记事本工具、Notepad++,
linux下的Vim、Emacs, 多端通用的 VSCode等；

编译器指将C、C++等程序代码编译为二进制可运行文件的程序, 常见的如随Visual Studio安装而安装的MSVC, linux下默认的GCC,
llvm组织出品的Clang等. 另一个常见的与之更类似的概念为解释器, 其为Python一类解释型语言所使用, 特点是从脚本逐行解释代码内容, 而不似编译器,
读取整个文件后编译为一个二进制可执行程序.；

IDE 为 Integrated Development Environment的缩写, 即集成开发环境, 可以看作是将编辑器和编译器结合起来的产物.
著名的Visual Studio即为集成开发环境, 此外还有JetBrains 公司出品的 CLion.

目前, MInDes使用Visual Studio开发Windows端, 使用VSCode + GCC + GDB + CMake的组合开发Linux端程序.

\subsection{文档(Documentation), API Reference}
程序开发领域的文档, 一般指帮助用户使用软件的用户手册, 亦或者记录软件功能的记录用手册. 在使用软件的过程中, 遇到的大多数问题都可以通过查询文档解决.
文档的书写风格通常不如教程来的简单易懂, 但文档会力求记录内容的准确性. 因此, 文档的编写某种程度也是程序开发不可缺少的一环, 而拥有较好的查阅文档的
能力也有助于使用软件或开发程序.

API Reference(参考)是与文档类似的概念, 所谓API指Application Programming Interface, 因此API Reference旨在向读者介绍程序开发中每一个函数/方法的
调用方式, 以助读者使用需要的API完成编程工作. 与文档相比, 更倾向于是给开发者使用的手册.

\subsection{编程语言, Python \& C++}
如今已经有众多的编程语言可供选择. 不同的编程语言有着不同的使用方法, 语言特性, 语法特点等. 目前常见的编程语言有Python, C, C++, Java,
C\#等等. 此处选择介绍Python和C++.

Python是一种面向对象的高级程序语言, 其语法灵活简单, 且因为丰富的三方库生态而充满活力. Python程序常由编写好的Python脚本(后缀常为\emph{.py})
开始, 经由python解释器进行逐行运行后给出程序结果, 因此也易于修改程序中的错误(Debug), 且现已有很多成熟的交互式Python编程方式, 如Jupyter等,
便于进行科学研究中的数据处理等. 但由于其内部函数封装层次太高, 运行效率通常不如更贴近机器层的C, C++等.

C++曾经作为C语言的面向对象版本出现, 而经过长期的开发已与C语言有了极大的区别. 然而C++仍然保留了许多C语言的特性, 和C语言在许多方面依旧是兼容的,
C++语法相比起Python较为复杂, 且调试过程也较为复杂, 但其运行效率更高, 更适合大规模的数据运算, 因而MInDes选择使用C++进行开发.

\subsection{面向对象}
面向对象是一种编程思想, 是相对于面向函数、面向过程等概念而生的, 其基本思想是一切皆为"对象", 对象内拥有必要的数据以及可以操作对象的方法,
且可以自由创建, 销毁对象, 通过调用相关的方法操作对象以实现程序的目的. 相比于顺序式执行的面向过程编程思想,
或是"一切皆函数"的面向函数思想, 面向对象有着更加灵活, 层次分明的特点, 适合于开发大型程序. 在C++中面向对象通过类(class)实现.

\subsection{C++语言标准}
C++ 历经长时间发展, 出现了许许多多不同的语言版本. 不同的语言版本之间的主要差别为语言特性的支持以及一些库是否可用.
目前最新的且较为广泛接受的语言版本为C++20, 普遍应用的版本则为C++14,
C++14标准同时也是Visual Studio 2022的默认语言标准, 而常称的"现代C++"则一般指代C++11及以后的标准, 因为C++11中引入了
大量的更现代的语法特性.

MInDes 出于兼顾稳定性以及功能性的原因, 选择使用C++17标准进行开发.

\subsection{编译}
编译(此处仅指C++的编译)是指将C++源代码通过编译器的处理生成一个可执行的程序的过程. 编译过程随着程序规模, 复杂程度而有时间上的变化, 但其基本过程可以
描述为以下几部分:
\begin{enumerate}
    \item 将源文件编译为对象文件(\emph{.o})
    \item 将对象文件按照源文件中所述的方式链接起来(此步骤需要链接器 (linker))
    \item 将链接后的得到的中间文件再与外部的静态/动态链接库相连接(此步也需要链接器)
    \item 生成二进制可执行文件
\end{enumerate}
从上述过程可知, 没有正确配置运行库环境可导致程序编译过程中产生链接器错误, 例如找不到符号的定义等. 因此开发环境的配置除了工具的设置外, 配置正确的运行
库也是环境配置的一部分.

\subsection{编译器, MSVC, GCC 与 C++标准委员会}
前面已经介绍了编译器的作用, 也介绍了C++语言标准. 语言标准由C++标准委员会商议确定, 通过对提案的可行性审查, 必要性审查等来确定标准C++在某个标准下应该
具有的语法特性. 该特性是推荐所有编译器共同支持的, 但因为各种原因, 部分语法特性并不是被所有编译器所支持的, 且每种不同的编译器也有各自不同的特点.
编译器可以实现非标准的C++语法特性, 为使用者提供便利, 但应注意使用了某款编译器特性的源代码, 在使用别的编译器时可能出现效率下降, 运行结果不符合预期,
甚至编译失败等各种问题. 其中, 处于Windows, Linux两大平台垄断地位的编译器, MSVC 与 GCC, 均拥有大量的为适应各自平台而开发的非标准的语法特性. 在跨平台
编程时, 应尽可能避免使用这类特性, 保持程序在双端都能有相似的表现.

MSVC编译器通常不直接接受编译器选项(flags / Compiler Options), 而是需要在VS中进行设置. 这类内容请查看Microsoft的各类文档, 如关于编译器选项的
\bhref{https://learn.microsoft.com/en-us/cpp/build/reference/compiler-options?view=msvc-170}{\underline{文档}}.

GCC是GNU Compiler Collection的缩写, 即GNU编译器套装. 历史上也代表GNU C Compiler以及GNU C \& C++ Compiler的缩写. GCC中与C++有关的组件为g++,
其可以通过在命令行后添加编译器指令(flags)以控制编译器的行为, 也可以通过各类构筑系统(下文将提到)以控制编译器的行为.

MSVC编译器的语法与GCC的有很大区别, 但一般不会特意手动设置这些内容.

\subsection{代码构筑系统, CMake}
一般的小型程序可以通过编译器简单的处理即可编译为可执行的二进制程序, 但当代码更加复杂, 文件众多且相互关联之时, 有一个辅助进行编译的程序可以有效地提高
程序开发效率. 这类辅助程序即代码构筑系统. 目前常见的构筑系统有Microsoft VS所采用的MSBuild构筑系统, Linux系统自带的Makefile, 广受欢迎的CMake,
图形化界面Qt所采用的QMake等等. 代码构筑系统为管理大型程序的编译, 调试等提供了便利. 目前MInDes采用MSBuild(Windows)与CMake(Linux)两种方式进行构建.

CMake的使用已经高度集成在VSCode中, 通过CMake插件提供的GUI可以简单实现CMake的各种操作. 然而, 依旧可以使用命令行进行操作. CMake的使用离不开
\emph{CMakeLists.txt}的编写, 通过插件可以实现基础的文件写入, 但为了实现更复杂的功能, 仍然需要自行编辑\emph{CMakeLists.txt}. 此处不再详细介绍
CMake 以及 \emph{CMakeLists.txt}的使用, 感兴趣可以查看CMake的文档, 但因为CMake为兼容并包多种写法, 文档也较为混乱, 建议必要时使用ChatGPT等AI
工具辅助生成\emph{CMakeLists.txt}.

\subsection{代码调试(Debug)}
代码一次写成, 运行良好固然很好, 但这种情况在实际开发中很难遇到. 实际开发中常会遇到各种各样的问题阻碍开发进展. 这些程序中或逻辑或语法的错误
就被称为Bug, 在程序中排查Bug并修正以使程序得以正常运行的过程即是调试, 亦即所谓的Debug.

Debug这一过程所占用的开发时间极长, 据统计几乎占据了实际开发用时的50\%以上, 因此拥有一定的Debug技巧将极大提高工作效率, 某种程度上也可以降低开发者的
心智负担. 最简单的调试方法即将程序在某一步的数据通过\cverb|cout|或\cverb|printf|输出到控制台上. 然而该方法的问题在于每次调试都需要改变代码,
重新编译, 然后再等待程序运行至对应的位置再将内容输出. 该方法具有多种局限性, 并不具备足够的实用性.

更现代化的技巧则为使用调试器(Debuger)运行程序
来实现类似于\\Python逐行解释运行的效果. VS中自带的MSVC已有调试器集成在内部, 而\\Linux平台下最常用到的Debuger即为gdb(GNU Debuger). 在调试器的辅助下,
可以在源代码中添加"断点"(一般表现为代码左端的红色圆点), 当程序在调试模式下编译后交由调试器启动时(通常IDE或CMake会将该步骤简化为一步), 程序运行至
断点所在行时便会停下, 不执行该行转而等待用户的命令. 此时可以通过调试控制台(Debug Console)和Debuger交互, 或者查看程序中变量的值, 甚至修改变量的值.
同时可以查看函数的调用栈情况, 也可以使用"步进"(Step Over)逐行运行, 或者"步入"(Step Into)查看该行函数内部的运行情况, 亦或者"步出"(Step Out)跳出该
函数回到上一层函数调用栈. Debuger的运用使程序的调试过程变得更加便利.

\subsection{并行运算, OpenMP与MPI}
并行运算是指通过将程序的运算工作按照某种方式拆分后分配给不同的算力单元进行运算, 以此充分提高运算效率, 减少资源浪费. 目前常见的并行运算方式有
OpenMP(Open Multi Process)这样多线程共享内存的并行方式, 也有MPI(Message Passing Interface)这样的多进程共享内存的并行方式. 

进程指代一个正在运行中的程序, 是系统资源调度分配的基本单位, 每个进程拥有自己的内存空间(即自己拥有独立的空间用以存储数据), 
而线程是进程的子任务, 是CPU调度分配资源的基本单位, 其并不单独占有某个内存空间, 而是从进程中获取自己所需的内存资源. 通过共享内存, 多个线程或进程之间
即可交换数据, 比如将一个大的数据划分为若干份, 每一份数据由一个线程(OpenMP)或一个进程(MPI)持有, 各个线程或进程之间有一定的数据交换以保证数据的
完整性. 在经过每个线程或进程对数据处理后, 再将处理好的结果汇总, 按照初始划分方式重新拼合起来. 线程所共享的内存空间只能是某一整个完整内存空间(所在的
进程的内存空间), 并不能跨内存空间交换数据, 因而OpenMP仅能支持单个内存空间内进行并行运算. 而MPI由于设计之初即为在多个进程之间交换数据, 亦即在多个
独立的内存空间之间交换数据, 绕过了线程不能跨内存空间的局限性. 目前大部分超算都支持MPI并行的程序.

OpenMP相比于MPI, 速度和可运算规模上MPI更加强大, 但MPI为实现多进程, 程序内部逻辑经常晦涩难懂, 且对程序的调试相较于OpenMP程序而言会复杂很多, 需要设置一个
主进程后令Debuger跟踪该主进程, 极大提高了调试的门槛, 且并不方便, 而OpenMP只需要在编译过程中给予编译器相关的设置, 便可以直接开始并行运算而无需复杂的操作. 
出于这样的考虑, 目前MInDes使用了OpenMP进行并行运算. 但未来可能考虑使用MPI以进一步提高软件的运算效率.

\subsection{Linux 内核与Linux 系统}
严格来讲, Linux应该指一种由Linus Torvalds所开发的操作系统内核, 意在模仿当时由贝尔实验室
开发的Unix操作系统, 实现一个完全自由的类Unix系统. 如今也被用以指代使用Linux内核的操作系统大类. 由于Linux内核仅作为操作系统的内核,
不同的组织或个人开发了不同的基于该内核的不同操作系统, 即为所称的Linux发行版. 目前有众多特性各异, 适用范围不同的发行版, 常见的个人
发行版为Ubuntu系统, 该发行版也是MInDes的主要测试目标系统.

\subsection{操作系统, WSL \& 虚拟机}
本程序的开发主力环境为Windows平台, 目前在Windows10 与 Windows11 平台上进行开发. 为使程序可以运行在算力更高的服务器或超算上,
本程序兼容了主流的Linux平台. 目前测试过的Linux发行版有Ubuntu与Arch Linux. 为了在Windows机器上可以测试Linux环境下程序运行情况,
一般的做法为使用虚拟机或者WSL（Windows Subsystem Linux）模拟出Linux操作系统环境.

其中, 虚拟机是一种成熟的技术, 通过在宿主机上的虚拟机软件内安装一个完整的Linux发行版来获得Linux环境. 而WSL则是微软所开发的一种Linux子系统,
和虚拟机相比, WSL的独立性不足, 但大部分文件是可以简单地共享的. 对本程序而言, WSL足以胜任大部分的Linux环境测试工作, 因而推荐使用WSL.

\subsection{命令行, 终端, Shell, Powershell \& bash}
开发程序的过程中, 尤其是在操作Linux系统时, 经常绕不开命令行操作. 命令行界面(Command Line Interface, CLI)是一种操作电脑的界面, 通过输入
文字命令以操作电脑执行命令, 通常程序初学者运行程序后出现的"黑框"即为一种CLI. 与之相对的概念即为图形用户界面(Graphical User Interface,
GUI). MInDes程序的操作建立在CLI上.

CLI需要有载体软件呈现, 这类载体被称为终端(Terminal). 最早的终端是一类实体机器, 通过终端输入命令来与计算机交互. 如今终端已经通过终端模拟器
内置于电脑中, 成为了一个软件. 与终端相关的概念为"控制台"(Console), 历史上二者有一定的区别, 现在大多不再区分.

Shell是为用户提供操作电脑方法的程序,在计算机层次架构中处于用户交互与底层运算之间的位置. 在不同的平台上均有众多各自不同的Shell程序, Windows端如
传统的cmd.exe, 较为现代化的Windows Powershell, 以及微软最新推出的Powershell7, 而Linux端则有常见的随系统安装自带的bash, 更加现代化, 自定义
程度更高, 与bash相兼容的zsh, 拥有与bash脚本不同语法, 更加人性化的fish等等.

Shell与命令行界面、终端之间的关系是: 终端内可运行多种Shell程序, Shell程序呈现出的界面为命令行界面. 在一般交流中三者区别并不明显,
影响交互方式的主要为Shell程序, 终端主要作为不同Shell程序的载体. \\在Windows端建议使用微软最新推出的Windows Terminal, 并且在Shell程序中推荐
命令与Linux有相似之处的Powershell7. 而在Linux端, 已有的bash已经具有较为强大的功能, 足以应对绝大部分场景的使用.
\endinput